\section{Diseño de experimentos}

\subsection{Escenarios}

\begin{table}[H]
\centering
\begin{tabular}{llc}
\toprule
\textbf{Escenario} & \textbf{Descripción} & \textbf{Enfermeras de triage} \\
\midrule
Base        & Configuración actual del servicio & 1 \\
Alternativo & Segundo puesto de clasificación en paralelo & 2 \\
\bottomrule
\end{tabular}
\caption{Escenarios comparados.}
\end{table}

El escenario alternativo modifica un único parámetro del modelo. La
estructura, las distribuciones y el resto de los supuestos permanecen
idénticos, lo que garantiza que cualquier diferencia observada sea atribuible
al cambio de dotación.

La justificación de este escenario es doble. Por una parte, el triage es la
única estación del sistema con un servidor único, configuración que en teoría
de colas resulta especialmente sensible a la variabilidad. Por otra, la cola
previa al triage es una cola ciega: todo minuto de espera anterior a la
clasificación transcurre sin priorización, de modo que reducirla beneficia
desproporcionadamente a los pacientes cuya urgencia todavía no fue
reconocida. Se trata además de la intervención de menor costo que el servicio
podría evaluar, frente a la alternativa de habilitar un quinto consultorio.

\subsection{Hipótesis de trabajo}

Se espera que la incorporación del segundo puesto reduzca de manera clara la
espera hasta la clasificación. Sin embargo, la mejora en el tiempo total de
permanencia debería resultar menor que el tiempo de espera de triage
eliminado, por dos razones: en régimen, el caudal de pacientes que alcanza la
cola de consultorios lo determinan los arribos y no el triage; y la
eliminación del puesto saturado suprime un efecto regulador que espaciaba la
entrada a la segunda cola.

La pregunta que responde la simulación es, en consecuencia, si el cuello de
botella se resuelve o simplemente se traslada.

\subsection{Horizonte y ventana de medición}

\begin{table}[H]
\centering
\begin{tabular}{ll}
\toprule
\textbf{Aspecto} & \textbf{Definición} \\
\midrule
Horizonte de simulación       & 300 minutos simulados (480 minutos reales) \\
Factor de escala temporal     & $k = 1{,}6$ \\
Período de calentamiento      & 45 minutos reales, descartados \\
Ventana de medición           & Minutos reales 45 a 405 \\
Réplicas por escenario        & 30 \\
Generación de números aleatorios & Semilla aleatoria por réplica \\
\bottomrule
\end{tabular}
\caption{Parámetros del diseño experimental.}
\end{table}

El período de calentamiento descarta el arranque con la guardia vacía, que no
representa una condición real de operación. El cierre anticipado de la
ventana de medición evita que los pacientes que no alcanzan a completar su
recorrido distorsionen los promedios.

La elección de treinta réplicas responde a la necesidad de disponer de una
muestra suficiente para estimar intervalos de confianza y aplicar el test de
medias. Con una única corrida por escenario la variabilidad entre réplicas
de una misma configuración puede superar la diferencia entre configuraciones,
de modo que ninguna conclusión sería sostenible.

\subsection{Variables de respuesta}

\begin{table}[H]
\centering
\small
\begin{tabular}{llp{6.2cm}}
\toprule
\textbf{Indicador} & \textbf{Unidad} & \textbf{Fundamento} \\
\midrule
Espera hasta el triage & min & Mide la cola ciega, sin priorización. \\
Espera hasta el consultorio & min & Detecta el traslado del cuello de botella. \\
Espera de pacientes verdes & min & Métrica~1 comprometida. \\
Espera de urgencias reales & min & Métrica adicional comprometida. \\
Permanencia total & min & Efecto neto sobre el paciente. \\
Pacientes en sala de espera & pacientes & Métrica~2, hacinamiento físico. \\
Utilización de los recursos & --- & Verificación y diagnóstico. \\
Recuento de sub y sobre-triage & pacientes & Calidad de la clasificación. \\
\bottomrule
\end{tabular}
\caption{Indicadores medidos y su justificación.}
\end{table}

La separación entre la espera de triage y la espera de consultorio resulta
esencial: si se midiera únicamente el tiempo total, un traslado del cuello de
botella sería indistinguible de la ausencia de efecto.

